\documentclass{situationist-preprint}

\title{The Lived Body as Zero-Point of Orientation}
\msid{PHE-08-a}
\msfamily{PHE-08}
\msdiscipline{Phenomenology}
\mstype{Principal manuscript}
\mskeywords{lived body; Leib; K\"orper; zero-point; orientation; double sensation; Husserl; Merleau-Ponty; reversibility; bodily awareness}

\begin{document}
\maketitle

\begin{abstract}
Is the lived body an object? Husserl distinguished the body as lived (\emph{Leib}) from the body as a physical thing (\emph{K\"orper}), and held that the lived body is the zero-point of orientation, the \enquote{here} from which everything else is there. Against this, it is often objected that my body is plainly something I can see, touch and inspect like any other thing, and that in double sensation, when one hand touches the other, the body is given to itself as an object without remainder. We argue that the lived body is not an object. We show that the zero-point of orientation cannot be among the things it orients, since an object is given as located relative to a here, and the here is not located relative to itself. We then analyse double sensation and show that it reveals an asymmetry rather than a symmetry: at any moment one hand is touching and the other touched, and the reversal of roles never coincides. Merleau-Ponty's thesis that reversibility is always imminent and never realized is the correct description. The body can be given as an object in part, but not fully, and the part that escapes is the part that does the orienting.
\end{abstract}

\section{The problem}

I can look at my hand, measure it, feel its warmth with the other hand, and see it in a mirror. In these respects my body is given to me as a thing among things, located in space, with a shape and a weight. Yet I do not stand to my body as I stand to a cup. I cannot walk around it, I cannot put it down, and I do not perceive it from a point of view, since it is the point of view from which I perceive. Husserl marked the difference by distinguishing \emph{Leib}, the lived body, from \emph{K\"orper}, the body as a physical thing, and held that the lived body is constituted in a way unlike any thing: it is the bearer of localized sensations, the organ of the will, and the zero-point of orientation \citep[\S\S18, 36--42]{Husserl1952}. Husserl also called it a \enquote{remarkably imperfectly constituted thing}, since parts of it, such as the back and the eyes, can never be seen by me directly \citep[\S41]{Husserl1952}.

Merleau-Ponty developed these themes into the thesis that the body is not an object but our means of having a world: it is never fully in front of me, it is always with me, and its permanence is not that of an object but the permanence of the perspective from which objects appear \citep{MerleauPonty1945}. Sartre distinguished the body as I exist it from the body as known by others, and held that the first is not an object for me at all \citep[Part~III, ch.~2]{Sartre1943}.

The opposing view has two forms. The first is a deflationary naturalism: the body is a physical object, and the differences Husserl describes are differences in how one kind of object is represented, through proprioception and interoception, not differences in what it is \citep{Bermudez1998,deVignemont2018}. The second appeals to double sensation. When my right hand touches my left, the left hand is felt as a thing, smooth, cool and located, and at the same time as sensing. Here, it is said, the lived body is given to itself as an object, and the distinction between \emph{Leib} and \emph{K\"orper} collapses into two aspects of one thing perceived.

We argue that the lived body is not an object, in the sense that it cannot be fully given to itself as an object. The argument has two parts: a structural argument from the zero-point (Section~\ref{s:zero}) and an analysis of double sensation (Section~\ref{s:double}).

\section{Framework}

\begin{definition}[Object]\label{d:obj}
Something is given as an \emph{object} to a subject if it is given as located in a space of positions relative to the subject's orientation, as having sides some of which are presented and some not, and as capable of being explored by a change in the subject's position.
\end{definition}

This is the sense of object in the perceptual tradition from Husserl on: a thing given in adumbrations, from a perspective, with a hidden side \citep[\S\S18, 41]{Husserl1952}. PHE-03 argues that every thing is given in adumbrations whose series does not converge on a complete givenness. An object in this sense is not merely a bearer of properties; it is something that appears as positioned relative to the one to whom it appears.

\begin{definition}[Zero-point of orientation]\label{d:zero}
The \emph{zero-point of orientation} of a subject at a time is that relative to which, at that time, everything perceived is given as near or far, left or right, above or below: the \enquote{here} of the subject's perceptual field.
\end{definition}

\begin{definition}[Full givenness as an object]\label{d:full}
Something is \emph{fully given to itself as an object} if, in some experience of a subject, it is given as an object to that subject in the sense of Definition~\ref{d:obj}, and nothing of it is at that time functioning otherwise than as given.
\end{definition}

\begin{principle}[Orientation is from the body]\label{p:orient}
For an embodied perceiver, the zero-point of orientation is the lived body: the here of the perceptual field is where the perceiving body is, and the directions of the field are fixed by the body's posture and the orientation of its organs.
\end{principle}

Principle~\ref{p:orient} is Husserl's thesis that the body is the bearer of the zero-point \citep[\S18]{Husserl1952}. It has independent support in the account of egocentric space: perceptual demonstratives such as \enquote{there} and \enquote{to the left} are understood through their connection to the perceiver's dispositions to act, and those dispositions are dispositions of a body \citep{Evans1982}. If I turn my head, what was in front is now to the side; the field's directions move with the body, not the body with them.

\section{The zero-point}\label{s:zero}

\begin{lemma}\label{l:notlocated}
The zero-point of orientation is not given, at the time at which it orients, as located in the space it orients.
\end{lemma}

\begin{proof}
By Definition~\ref{d:obj}, to be given as located is to be given as at some position relative to the subject's here. By Definition~\ref{d:zero}, the zero-point is the here. If it were given as located relative to the here, it would be given as at some distance and direction from itself. But the here is at no distance and in no direction from itself: it is that relative to which distances and directions are fixed. It therefore cannot be given as located in the field that it orients. It can be given in other ways, as felt from within or as the place from which one acts, but these are not the ways an object is given.
\end{proof}

\begin{theorem}\label{t:notobj}
The lived body is not fully given to itself as an object.
\end{theorem}

\begin{proof}
Consider any experience in which the lived body is given to the subject. By Principle~\ref{p:orient}, in that experience the lived body is functioning as the zero-point of orientation, since the subject is perceiving and the here of the field is the perceiving body. By Lemma~\ref{l:notlocated}, the body so functioning is not given as located in the field it orients. Hence something of the body is, at that time, functioning otherwise than as given, namely as the zero-point. By Definition~\ref{d:full}, the body is not fully given to itself as an object in that experience. Since the experience was arbitrary, it is not so given in any experience.
\end{proof}

The theorem does not say that no part of my body can be given as an object to me. When I look at my hand, the hand is given as located: there, in front of me, below my eyes. But the eyes and the head that orient the visual field are not, at that moment, given as located in it. One may then look at the eyes in a mirror; but the mirror image is given as there, and what orients the field is still here, now including the eyes that see the image. Every attempt to make the orienting body an object displaces the orientation to another part of the body or to the same part functioning differently. Husserl's remark about the imperfectly constituted thing is a symptom of this structure: what cannot be seen directly is precisely what does the seeing.

\section{Double sensation}\label{s:double}

The strongest case for the objectivity of the body is double sensation (\emph{Doppelempfindung}), which Husserl analyses in \S\S36--37 of \emph{Ideas~II}. When my right hand touches my left, two series of sensations arise. In the touching hand there are sensations that present the touched hand as a thing, smooth, cool, with a certain shape. In the touched hand there are sensations localized in it, felt as \emph{in} it: the pressure is felt at the place where it is touched. Husserl argues that it is through such localized sensations, which he calls \emph{Empfindnisse}, that the body is constituted as lived; a merely seen body would not be lived in this way \citep[\S\S36--37]{Husserl1952}.

The objection runs: in double sensation the left hand is both lived (bearer of sensations) and a thing (felt as smooth and cool by the right), and the right hand can at once be felt as a thing by the left. So the body is given to itself both as subject and as object, without remainder.

\begin{definition}[Roles]\label{d:roles}
In a tactile experience in which one part $A$ of the body touches another part $B$, $A$ plays the \emph{touching role} if through $A$'s sensations $B$ is given as an object (Definition~\ref{d:obj}); $B$ plays the \emph{touched role} if $B$ is so given.
\end{definition}

\begin{proposition}[Asymmetry]\label{p:asym}
At any moment of double sensation, the touching and touched roles are not both played by each hand with respect to the other. The roles can be reversed, but the reversal is a succession, not a coincidence.
\end{proposition}

\begin{proof}
Suppose that at a moment $t$ the right hand plays the touching role with respect to the left and the left plays the touching role with respect to the right. Then each hand is given as an object through the other's sensations, so each is given as located relative to the subject's here and as having presented and hidden sides. The tactile field at $t$ would then have two zero-points, since the exploring hand is the here of a tactile exploration (Principle~\ref{p:orient}): the right hand would be the here relative to which the left is there, and the left the here relative to which the right is there. By Lemma~\ref{l:notlocated} the here of a field is not located in it. So the right hand would be both located in the tactile field, as touched, and not located in it, as its here. This is impossible. Hence at $t$ at most one hand plays the touching role. That the roles can be reversed is shown by experience: attending to the left hand's exploration makes it the touching hand and the right the touched one. But the reversal displaces the here; it does not make both hands the here at once.
\end{proof}

Proposition~\ref{p:asym} is Merleau-Ponty's description. The two hands are never simultaneously touched and touching with respect to each other; when I press them together, I do not have two sensations felt together as one perceives two objects placed side by side, but an ambiguous arrangement in which the two hands can alternate the roles of touching and being touched \citep{MerleauPonty1945}. In his last work he made reversibility the structure of the flesh, adding that it is always imminent and never realized in fact: either my left hand is touching the right, or my right hand feels the left, but I never reach the coincidence \citep{MerleauPonty1964}. The phenomenon that was supposed to show that the body is given to itself fully as an object shows instead that it never is: in every moment of the experience something of the body is not given but giving.

The asymmetry also answers the deflationary naturalist. \citet{Bermudez1998} argues that bodily awareness is a form of perception of the body as an object, distinctive in its sources but not in its object, and \citet{deVignemont2018} that the sense of bodily ownership is the content of a representation of a protective body map. These are accounts of how the body is represented, and they may be true. They do not show that the body is fully given as an object, because they describe the body as represented, which is the touched role; the touching or perceiving role is the representing system's own functioning, which is not in the same experience represented as located. \citet{Martin1995} argued that in bodily sensation one is aware of a region as falling within the boundaries of one's body, and that this awareness has a spatial structure; but the space so structured is the body's own felt space, in which locations are given relative to the body's posture and not relative to an external here. That is the space of the lived body, not the space of objects.

\section{Objections and replies}

\begin{objection}[Third-person givenness]\label{o:third}
Others see my body as an object, and so can I in a photograph. What is given to others as an object is an object.
\end{objection}

\begin{reply}
The thesis concerns givenness to oneself. My body is an object for others, and through their look I can grasp it as such, as Sartre's third ontological dimension of the body records \citep[Part~III, ch.~2]{Sartre1943}. The photograph gives a body there, which I identify as mine. Neither experience is one in which the body that is perceiving is given as perceived. The question was whether the lived body can be fully given \emph{to itself} as an object, and neither case is a case of that.
\end{reply}

\begin{objection}[Successive totality]\label{o:succ}
Even if no single experience gives the body fully as an object, a series of experiences can, by reversals and mirrors, give every part of it as an object in turn.
\end{objection}

\begin{reply}
Each experience in the series gives some parts as objects and leaves others functioning. The union of what is given in the series may include every part, but no experience gives the whole as object, and the parts are not given together as one located thing seen from a here, since the here moves with each step. Moreover, the series does not converge: each step leaves a functioning remainder, as every perspectival series leaves a hidden side (PHE-03). The body is given as an object in pieces and in succession, which is not to be given fully as an object.
\end{reply}

\begin{objection}[Proprioceptive location]\label{o:prop}
I know where my limbs are without looking, and know it in a spatial frame shared with external objects: I can point to my left knee with my eyes closed. So my body is given as located.
\end{objection}

\begin{reply}
Proprioceptive location is location relative to the body's own axes, the felt space of Martin's account, integrated with external space through the body schema \citep{Gallagher2005}. Gallagher distinguishes the body schema, the system of sensorimotor capacities that functions without perceptual monitoring, from the body image, the system of perceptions and beliefs about one's body. The knee is given as located through the body image; the body schema that orients the pointing is not itself given as located in that act. The objection therefore confirms the structure: part of the body is given as there, while the body as orienting functions as here.
\end{reply}

\begin{objection}[Merely epistemic]\label{o:epis}
The asymmetry shows only a limitation of perspective, not that the lived body is not an object. A camera cannot photograph its own lens, but the lens is an object.
\end{objection}

\begin{reply}
The analogy concedes the point. The lens is an object for someone else, and could be photographed by another camera. The question is whether the lived body is an object in the order of its own givenness, and the camera case shows only that a viewing system cannot view its own viewing. For the body the limitation is not accidental: there is no other perceiver that is my perceiving body, so no standpoint from which it could be given to me as the lens can be given to another camera. The thesis is therefore not that the body lacks objective properties, but that its being as lived, the being in virtue of which it is the zero-point, is not the being of an object. \citet{Zahavi2005} and \citet{Legrand2006} have argued that pre-reflective bodily self-awareness is a form of self-awareness of the body as subject, not as object; the asymmetry of double sensation is its clearest phenomenological mark.
\end{reply}

\section{What the result does not establish}

The result does not show that the body has no physical properties, or that it is not a physical object in the sense of physics; it concerns givenness. It does not establish any form of dualism, since the same body can be lived and physical. It does not settle whether the lived body is the self, or whether there can be persons without bodies; it shows only that for an embodied perceiver the zero-point is bodily. And it does not establish that the body's givenness as lived is incorrigible: one can misjudge the position of a limb, as in phantom-limb cases, without the zero-point ceasing to function as such.

\section{Conclusion}

The lived body is not an object. It is the zero-point of orientation, and the zero-point cannot be given as located in the field it orients. Parts of the body can be given as objects, and the whole can be given as an object to others, but the body is never fully given to itself as an object, because something of it always functions as the here. Double sensation, the case most often cited against this view, confirms it: the two hands can exchange the roles of touching and touched, but at no moment does each play both, and the coincidence of the lived and the perceived body is always imminent and never realized.

\begin{thebibliography}{99}\small
\bibitem[Berm\'udez(1998)]{Bermudez1998} Berm\'udez, J.~L. (1998). \emph{The Paradox of Self-Consciousness}. MIT Press.
\bibitem[de~Vignemont(2018)]{deVignemont2018} de Vignemont, F. (2018). \emph{Mind the Body: An Exploration of Bodily Self-Awareness}. Oxford University Press.
\bibitem[Evans(1982)]{Evans1982} Evans, G. (1982). \emph{The Varieties of Reference}. Ed. J. McDowell. Oxford University Press.
\bibitem[Gallagher(2005)]{Gallagher2005} Gallagher, S. (2005). \emph{How the Body Shapes the Mind}. Oxford University Press.
\bibitem[Husserl(1952)]{Husserl1952} Husserl, E. (1952). \emph{Ideen zu einer reinen Ph\"anomenologie und ph\"anomenologischen Philosophie. Zweites Buch}. Husserliana IV. Ed. M. Biemel. Martinus Nijhoff. Trans. R. Rojcewicz and A. Schuwer, \emph{Ideas Pertaining to a Pure Phenomenology and to a Phenomenological Philosophy, Second Book}, Kluwer, 1989.
\bibitem[Legrand(2006)]{Legrand2006} Legrand, D. (2006). The bodily self: The sensori-motor roots of pre-reflective self-consciousness. \emph{Phenomenology and the Cognitive Sciences}, 5(1), 89--118.
\bibitem[Martin(1995)]{Martin1995} Martin, M.~G.~F. (1995). Bodily awareness: A sense of ownership. In J.~L. Berm\'udez, A. Marcel and N. Eilan (eds.), \emph{The Body and the Self}. MIT Press.
\bibitem[Merleau-Ponty(1945)]{MerleauPonty1945} Merleau-Ponty, M. (1945). \emph{Ph\'enom\'enologie de la perception}. Gallimard. Trans. D.~A. Landes, \emph{Phenomenology of Perception}, Routledge, 2012.
\bibitem[Merleau-Ponty(1964)]{MerleauPonty1964} Merleau-Ponty, M. (1964). \emph{Le visible et l'invisible}. Ed. C. Lefort. Gallimard. Trans. A. Lingis, \emph{The Visible and the Invisible}, Northwestern University Press, 1968.
\bibitem[Sartre(1943)]{Sartre1943} Sartre, J.-P. (1943). \emph{L'\^etre et le n\'eant}. Gallimard. Trans. H.~E. Barnes, \emph{Being and Nothingness}, Philosophical Library, 1956.
\bibitem[Zahavi(2005)]{Zahavi2005} Zahavi, D. (2005). \emph{Subjectivity and Selfhood: Investigating the First-Person Perspective}. MIT Press.
\end{thebibliography}

\end{document}
